\documentclass[11pt]{article}
\usepackage[margin=1in]{geometry}
\usepackage{amsmath,amsthm,amssymb,hyperref,cleveref}
\newtheorem{theorem}{Theorem}[section]
\newtheorem{lemma}[theorem]{Lemma}
\newtheorem{proposition}[theorem]{Proposition}
\newtheorem{corollary}[theorem]{Corollary}
\begin{document}
\section{Elementary integral chain constructions}
\label{mapsI:chain:section}

These classical constructions are included as rederivations of the singular-chain prerequisites; they make no novelty claim.

\subsection{Chains, pairs and homotopies}
\label{mapsI:chain:section-1}

Let \(C_m(X)\) be the free abelian group on the continuous maps \(\sigma:\Delta^m\to X\), with \(C_m=0\) for \(m<0\). Write \(\sigma[i_0,\ldots,i_k]\) for its restriction to the indicated ordered face. The differential is

\[
\partial\sigma=\sum_{i=0}^m(-1)^i\sigma[0,\ldots,\widehat i,\ldots,m].
\]

Each twice-deleted face appears with opposite signs in \(\partial^2\sigma\), so \(\partial^2=0\). Homology is \(H_m=\ker\partial/\operatorname{im}\partial\). For \(A\subset X\), the relative complex is \(C_*(X,A)=C_*(X)/C_*(A)\). It is free on the simplices not entirely contained in \(A\). With an abelian coefficient group \(G\), use \(C_*(X,A)\otimes G\) for homology and \(\operatorname{Hom}(C_*(X,A),G)\) for cohomology; its coboundary is \(\delta a=a\partial\). Maps of pairs induce the evident chain maps and contravariant cochain maps.

\begin{lemma}[homotopy invariance]
\label{mapsI:chain:lemma:1.1}
Homotopic maps of pairs induce the same homology and cohomology maps, with every coefficient group.
\end{lemma}

\begin{proof}
For a homotopy \(H:X\times I\to Y\), triangulate \(\Delta^m\times I\) into the \(m+1\) ordered simplices

\[
[(v_0,0),\ldots,(v_i,0),(v_i,1),\ldots,(v_m,1)],
\qquad 0\leq i\leq m.
\]

Apply \(H(\sigma(-),-)\) to their sum with coefficients \((-1)^i\), defining \(P\sigma\). Adjacent internal faces cancel. The bottom and top faces remain with signs minus and plus, and the side faces are \(-P\partial\sigma\). Hence
\(\partial P+P\partial=H_{1*}-H_{0*}\).
If \(H(A\times I)\subset B\), this formula descends to relative chains. Tensoring the homotopy, or precomposing cochains with it, proves the assertion.
\end{proof}

The point complex has one generator in every nonnegative degree; its differential is zero in odd degrees and the identity in positive even degrees. Therefore a point has homology \(\mathbb Z\) in degree zero and zero elsewhere. By \cref{mapsI:chain:lemma:1.1} the same is true of a contractible space. Moreover \(H_0(X)\) is free on the path components: the boundary of a path equates its two endpoints, and these are exactly the relations imposed by paths.

\subsection{Making singular chains small}
\label{mapsI:chain:section-2}

For an open cover \(\mathcal U\) of \(X\), let \(C_*^{\mathcal U}(X)\) be generated by simplices whose image lies in one cover member. Local arguments require the inclusion of these chains to be a chain homotopy equivalence, rather than just an intuitive subdivision picture.

\begin{lemma}[small-chain lemma]
\label{mapsI:chain:lemma:2.1}
There is a chain map \(\rho:C_*(X)\to C_*^{\mathcal U}(X)\) and a homotopy \(D\) with

\[
\rho\iota=1,\qquad 1-\iota\rho=\partial D+D\partial.
\]

Both preserve every subspace of \(X\), so the assertion also holds for relative complexes.
\end{lemma}

\begin{proof}
First construct subdivision \(S\) on the affine chains in each standard simplex. If \(b\) is its barycentre, define \(S\sigma=b*(S\partial\sigma)\) inductively, starting with vertices. Here the affine cone with apex first satisfies
\(\partial(b*c)=c-b*(\partial c)\) for a positive-dimensional cycle, and the corresponding reduced formula for dimension zero. It follows inductively that \(\partial S=S\partial\). Composing these affine subdivisions with a singular simplex defines a natural chain map on singular chains.

Construct a subdivision homotopy \(T\) by induction too. On an \(m\)-simplex, the chain
\(z=\sigma-S\sigma-T\partial\sigma\) is a cycle by the already established lower-dimensional identity. It has augmentation zero when relevant. Cone it to the barycentre and set \(T\sigma=b*z\). Then
\(\partial T+T\partial=1-S\).
All chains constructed remain inside the original simplex image. Thus subdivision and its homotopy preserve small chains and chains in any subspace.

A barycentric subsimplex has vertices that are barycentres of nested faces. If a face of size \(\ell\) lies in one of size \(k\), their barycentre difference has norm at most \((k-\ell)/k\) times the old diameter, at most \(m/(m+1)\) times it. The diameter of a Euclidean simplex is the maximum of its vertex distances. Iteration therefore makes the mesh tend to zero. For the inverse-image open cover of the compact standard simplex, choose at each point a ball whose doubled ball lies in one cover member. Finitely many of the smaller balls cover the simplex; the minimum of their radii is positive. Every subset of smaller diameter, by choosing one of its points in a smaller ball, lies in the corresponding doubled ball. This is the required Lebesgue number. Every singular simplex, and every finite chain, becomes small after sufficiently many subdivisions.

Put \(T_N=\sum_{j=0}^{N-1}S^jT\); then
\(\partial T_N+T_N\partial=1-S^N\).
Now construct \(\rho,D\) inductively on the basis of singular simplices. Given their values on the boundary of \(\sigma\), put
\(z=\sigma-D\partial\sigma\), so \(\partial z=\rho\partial\sigma\) is small. Choose \(N\) such that \(S^Nz\) is small and define

\[
\rho\sigma=S^Nz+T_N(\rho\partial\sigma),
\qquad D\sigma=T_Nz.
\]

The second term of \(\rho\sigma\) is small because \(T_N\) preserves smallness. The displayed identities give
\(\partial\rho\sigma=\rho\partial\sigma\) and
\(\sigma-\rho\sigma=\partial D\sigma+D\partial\sigma\).
If \(\sigma\) is already small, its faces have \(D=0\) by induction; choose \(N=0\), obtaining \(\rho\sigma=\sigma\), \(D\sigma=0\). Every operation remains in the original image, which proves the subspace assertion.
\end{proof}

\begin{corollary}[excision]
\label{mapsI:chain:corollary:2.2}
If \(Z\subset A\subset X\) and \(\overline Z\subset\operatorname{int}_X A\), the inclusion
\((X\setminus Z,A\setminus Z)\to(X,A)\)
induces homology and cohomology isomorphisms with every coefficient group.
\end{corollary}

\begin{proof}
Use the open cover \(\operatorname{int}A, X\setminus\overline Z\) of \(X\). After quotienting small chains by those in \(A\), the first summand disappears and the quotient is
\(C_*(X\setminus\overline Z,A\setminus\overline Z)\).
The restricted cover of \(X\setminus Z\) gives the very same quotient: its first member is \((\operatorname{int}A)\setminus Z\), whose chains again disappear. \cref{mapsI:chain:lemma:2.1} identifies both ordinary relative complexes with these quotients up to chain homotopy. Tensor and dualize the equivalences.
\end{proof}

\subsection{Exact sequences and local gluing}
\label{mapsI:chain:section-3}

A degreewise exact sequence of complexes \(0\to K\to L\to M\to0\) has a long exact sequence in homology. Its connecting map sends a cycle \(m\) to the class of \(\partial\ell\) in \(K\), where \(\ell\) lifts \(m\). Changing the lift changes this by a boundary; changing the cycle by a boundary does likewise. Exactness follows directly: a cycle of \(L\) maps to a boundary of \(M\) exactly when, after subtracting the boundary of a lift, it lies in \(K\); a cycle of \(M\) has zero connecting class exactly when its lift can be corrected by a chain in \(K\) to a cycle of \(L\); a cycle of \(K\) becomes a boundary in \(L\) exactly when it is the connecting image of the projected bounding chain. These three arguments prove exactness at all the repeating positions. The same lifting argument, with \(\delta\), proves the cohomological version and its naturality.

The pair sequence comes from
\(0\to C_*(A)\to C_*(X)\to C_*(X,A)\to0\).
It is degreewise split on the simplex bases, so it stays exact after tensoring or dualizing.

For an open cover \(X=U\cup V\), the sequence

\[
0\to C_*(U\cap V)\xrightarrow{c\mapsto(c,-c)}
C_*(U)\oplus C_*(V)\xrightarrow{(a,b)\mapsto a+b}
C_*^{\{U,V\}}(X)\to0
\]

is degreewise split exact: each simplex belongs to \(U\), to \(V\), or to both. \cref{mapsI:chain:lemma:2.1} replaces the last complex by \(C_*(X)\). We obtain the \textbf{Mayer--Vietoris sequences}. For cohomology the middle-to-intersection map is the difference of restrictions. Dividing throughout by the corresponding chains in a subspace \(A\) gives the relative sequences for \((X,A)\), with pieces \((U,U\cap A)\), \((V,V\cap A)\) and \((U\cap V,U\cap V\cap A)\). Splitting is again obtained on the simplices outside \(A\).

We will use the following precise form of the five lemma. In a commutative diagram of exact five-term sequences, if maps at positions one, two, four and five are isomorphisms, the middle map is an isomorphism. For injectivity, an element with zero image maps to zero at position four, so comes from position two; lift its image-zero condition through position one and use the isomorphisms there to correct its preimage to zero. For surjectivity, map a target middle element to position four; lift there, note its image at position five vanishes, and lift back to a source middle element. Their difference lies in the image of position two, whose isomorphism provides the required correction. This proves the lemma and applies equally to homology and cohomology segments.

For example, cover \(S^m\), \(m\geq1\), by slightly enlarged hemispheres. Each is contractible, and their intersection retracts to \(S^{m-1}\). The reduced Mayer--Vietoris sequence gives
\(\widetilde H_j(S^m;G)\cong\widetilde H_{j-1}(S^{m-1};G)\).
Starting with two points, this yields \(\widetilde H_m(S^m;G)=G\) and zero in other degrees; the same argument gives reduced cohomology. The generator obtained by orienting the two hemispheres compatibly is the fundamental class of the oriented sphere. Concretely triangulate it as the boundary of an \(m+1\)-simplex, radially projected from an interior point. Each ray meets the boundary once, and its intersection distance is continuous face by face, so this is a homeomorphism. The alternating sum of oriented faces is a cycle because shared faces cancel. Its image at a point inside one face is the positive local disk generator by excision. Since the sphere homology group in degree \(m\) is infinite cyclic, that value one proves the cycle is a generator. To verify the same local value at any point, choose a different centre or a rotated simplex whose face interior contains that point; the radial projections are related through orientation-preserving rotations, which are homotopic to the identity by the following elementary plane-rotation argument: an orthogonal orientation-preserving matrix is reduced to the identity by successive rotations in coordinate two-planes, each connected to the identity by its angle parameter. This gives the fundamental class \([S^m]\), with local value one in every oriented chart.

\subsection{Products, signs and evaluation on fibres}
\label{mapsI:chain:section-4}

For cochains of degrees \(p,q\) with values in a commutative ring \(R\), define the Alexander--Whitney cup product by

\[
(a\smile b)(\sigma)=a(\sigma[0,\ldots,p])
b(\sigma[p,\ldots,p+q]).
\]

Expanding boundaries, terms deleting vertices strictly before the joining vertex give \(\delta a\smile b\), those after it give \((-1)^p a\smile\delta b\), and the two joining terms cancel. Thus

\[
\delta(a\smile b)=\delta a\smile b+(-1)^p a\smile\delta b.
\]

The formula is natural. It is associative: the same three face intervals occur in either order of multiplication. The constant zero-cochain \(1\) is its unit. It descends to cohomology. An absolute class can multiply a relative class for every pair; cochains on the subspace vanish in the relative factor.

We supply the chain comparison that explains external products and graded commutativity. On the total tensor complex use
\(\partial(c\otimes d)=\partial c\otimes d+(-1)^{\deg c}c\otimes\partial d\).
Define

\[
A\sigma=\sum_{p+q=m}\sigma_X[0,\ldots,p]\otimes\sigma_Y[p,\ldots,m],
\qquad \sigma:\Delta^m\to X\times Y.
\]

Conversely define \(S(c\otimes d)\) by triangulating \(\Delta^p\times\Delta^q\) along all lattice paths from \((0,0)\) to \((p,q)\). Give a path the sign of the permutation of its horizontal and vertical steps relative to all horizontal steps followed by all vertical steps. Compose its affine simplex with \((c,d)\). Internal faces from paths interchanging adjacent unlike steps cancel; boundary faces give precisely \(\partial c\otimes d+(-1)^p c\otimes\partial d\). For \(A\), deleting a nonjoining vertex gives the matching tensor boundary and joining-vertex terms telescope. Hence both maps are chain maps, preserving augmentation.

\begin{lemma}[product chain equivalence]
\label{mapsI:chain:lemma:4.1}
\leavevmode\par
The two maps are inverse up to natural chain homotopies, preserving products of subspaces.
\end{lemma}

\begin{proof}
Here is an explicit inductive construction of the needed homotopies. Chains in a convex set have a contraction to any fixed vertex: cone each singular simplex to that vertex by linear interpolation. The cone boundary identity gives \(\partial K+K\partial=1-i\epsilon\), where \(\epsilon\) is augmentation and \(i\) inserts the vertex. Tensoring two such contractions gives a contraction of their tensor complex to \(\mathbb Z\): the homotopy is \(K\otimes1+(i\epsilon)\otimes L\), with the tensor sign \((-1)^{\deg c}\) in the second term. The first identity and then the second verify this formula; its boundary sum is \(1-(i\epsilon)\otimes(j\epsilon)\).

For \(1-SA\), assume a homotopy \(h\) is already defined on lower-dimensional simplices. On the universal diagonal simplex \(\Delta^m\to\Delta^m\times\Delta^m\), form
\(z=(1-SA)\sigma-h\partial\sigma\).
It is an augmented cycle. Cone \(z\) in the convex product to fill it. For an arbitrary \(\sigma\), push this filling forward by \(\sigma_X\times\sigma_Y\). This defines \(h\) naturally and gives \(\partial h+h\partial=1-SA\).

For \(1-AS\), repeat on a universal tensor generator \(1_{\Delta^p}\otimes1_{\Delta^q}\). The residual cycle is filled by the tensor contraction just constructed, then pushed forward by the two singular simplices. Induction on \(p+q\) gives the other homotopy. If either simplex lies in a subspace, every filling still maps into the corresponding product subspace.
\end{proof}

The same induction compares \(A\) with the map obtained by swapping the two coordinates and then applying \(A\) and the signed tensor swap
\(c_p\otimes d_q\mapsto(-1)^{pq}d_q\otimes c_p\).
They agree on augmentation, and the tensor contraction fills the residual cycles exactly as in \cref{mapsI:chain:lemma:4.1}. Evaluate \(a\otimes b\) on this homotopy after the diagonal \(X\to X\times X\). For cocycles the resulting difference is a coboundary. Therefore

\[
[a]\smile[b]=(-1)^{pq}[b]\smile[a].
\]

No assertion of cochain commutativity was used.

The equivalence also descends to
\(C_*(X,A)\otimes C_*(Y,B)\) and the product complex modulo
\(C_*(A\times Y)+C_*(X\times B)\).
When \(A,B\) are open, small chains in their union identify this quotient up to chain homotopy with
\(C_*(X\times Y,A\times Y\cup X\times B)\).
If one subspace is empty the assertion holds for the other subspace without openness. For classes, define the external product by \(a\times b=\operatorname{pr}_1^*a\smile\operatorname{pr}_2^*b\); its cochain formula is evaluation of \(a\otimes b\) on \(A\). This defines external products on all pairs used here. Cup products of two relative classes likewise take values relative to their union when its two subspaces are open, by the same small-chain comparison.

For an integral relative cocycle \(u\) of degree \(r\), define a right cap operator on an \(m\)-simplex, \(m\geq r\), by

\[
R_u\sigma=u(\sigma[m-r,\ldots,m])\,\sigma[0,\ldots,m-r],
\]

and set it to zero when \(m<r\). Boundary expansion and \(\delta u=0\) give
\(\partial R_u=R_u\partial\); the faces at the cut cancel by the cocycle equation on the last \(r+1\) vertices. Thus it is a degree-\(-r\) map commuting with the unshifted differentials. If \(u\) vanishes on a subspace, this map descends from relative chains to absolute chains. For a cocycle \(a\) of degree \(m-r\), its defining identity is

\[
a(R_u\sigma)=(a\smile u)(\sigma).
\]

We use this convention throughout. Reindexing chain groups here keeps their differential unchanged; it is not the alternative signed convention for suspended chain complexes. Replacing \(u\) by \(u+\delta v\) changes the cap map by a homotopy: on \(C_m\) set \(h=(-1)^{m-r+1}R_v\). The cup-product differential identity gives \(\partial h+h\partial=R_{\delta v}\). Products and this cap identity also make sense with an abelian coefficient group in the first factor and an integral cochain in the second, using integer multiplication on the group.

\section{Actual CW models and the first Hurewicz map}
\label{mapsI:cw:section}

The singular-chain prerequisites, sphere generators and prism calculations are proved in \cref{mapsI:chain:section}. The first Hurewicz theorem is the classical result of \cite[Theorem~4.32]{Hatcher2002}; the following is our complete expanded proof.

\subsection{A CW model for any open Euclidean neighbourhood}
\label{mapsI:cw:cw-open}

Let \(U\subset\mathbb R^N\) be open. For \(n\geq1\), put \(h_n=2^{-n}\). Let \(K_n\) be the union of all closed grid cubes of side \(h_n\) contained in \([-n,n]^N\) and satisfying

\begin{equation}
\operatorname{dist}_{\infty}(Q,\mathbb R^N\setminus U)\geq h_n.
\label{mapsI:cw:eq:1.1}
\end{equation}

Include every face of every chosen cube. If the complement is empty, its distance is \(+\infty\). Thus \(K_n\) is a finite cubical complex contained in \(U\); empty early stages cause no problem.

Every child cube of an eligible cube is eligible at the next stage. Hence \(K_n\subset K_{n+1}\). In fact

\begin{equation}
K_n\subset\operatorname{int}K_{n+1},
\qquad
\bigcup_n\operatorname{int}K_n=U.
\label{mapsI:cw:eq:1.2}
\end{equation}

For the first assertion, the layer of next-grid cubes immediately surrounding an eligible \(Q\) is at sup-distance at most \(h_{n+1}\) from \(Q\). Its distance from the complement is therefore at least \(h_n-h_{n+1}=h_{n+1}\). It lies in the next bounding box, so that full surrounding layer is eligible. For the second assertion, a point of \(U\) has positive distance from the complement on a sufficiently small compact neighbourhood. Choose the grid fine enough that this neighbourhood's intersecting cubes satisfy \eqref{mapsI:cw:eq:1.1}, and the bounding box large enough to contain them. It then lies in \(\operatorname{int}K_n\).

Give each \(K_n\) its grid-cell structure. The inclusion \(K_n\hookrightarrow K_{n+1}\) is cellular: an old \(p\)-face is a union of next-grid \(p\)-faces, and all those faces belong to eligible child cubes. Form the actual mapping telescope

\begin{equation}
T=\left(\coprod_{n\geq1}K_n\times[n,n+1]\right)\big/\bigl((x,n+1)_n\sim(x,n+1)_{n+1}\bigr).
\label{mapsI:cw:eq:1.3}
\end{equation}

Here the identification uses the original inclusion of \(K_n\) into \(K_{n+1}\).

The telescope is a CW complex. On each stage retain the bottom and top grid cells and attach a prism for each cell of \(K_n\). A \(p\)-cell's prism is a \((p+1)\)-cell; its boundary consists of the old cell, its cellular image at the next level, and the already attached prisms of its boundary. The image at the next level contains only finitely many cells. Attaching in dimension order gives the required characteristic maps and closure finiteness. Every bounded height interval meets finitely many stages and cells; the resulting CW topology agrees with the telescope's quotient topology.

More concretely \(T\) is the following subspace of \(U\times[1,\infty)\):

\begin{equation}
T=\bigcup_{n\geq1}K_n\times[n,n+1].
\label{mapsI:cw:eq:1.4}
\end{equation}

The two models agree topologically, not just as sets. On a bounded height interval there are finitely many closed compact slabs. Their union has exactly the topology obtained by gluing their common closed subsets. Such intervals cover the telescope locally, proving the assertion. Let \(q:T\to U\) be projection to the first coordinate.

We construct a continuous section and homotopy explicitly. For each \(n\), choose a smooth function \(\chi_n:U\to[0,1]\), equal to one on a neighbourhood of \(K_n\), supported in \(\operatorname{int}K_{n+1}\). These exist by \eqref{mapsI:cw:eq:1.2}: cover the compact set \(K_n\) by finitely many Euclidean balls whose larger concentric balls lie in that interior, choose smooth radial bumps \(b_i\) equal to one on the smaller balls, and use \(1-\prod_i(1-b_i)\). Take \(\chi_n=0\) if \(K_n\) is empty. Define

\begin{equation}
\rho(x)=3+\sum_{n\geq1}(1-\chi_n(x)).
\label{mapsI:cw:eq:1.5}
\end{equation}

This sum is locally finite. On a neighbourhood contained in \(K_N\), every term with \(n\geq N\) is zero. If \(j\) is the first index for which \(x\in K_j\), then \(\chi_n(x)=0\) for \(n\leq j-2\). Hence \(\rho(x)\geq j+1\) when \(j\geq2\), and \(\rho(x)\geq3\) when \(j=1\). In either case \(x\in K_{\lfloor\rho(x)\rfloor}\). We have

\begin{equation}
s:U\longrightarrow T,\quad s(x)=(x,\rho(x)),\qquad q\,s=1_U.
\label{mapsI:cw:eq:1.6}
\end{equation}

Both maps are continuous in the subspace model \eqref{mapsI:cw:eq:1.4}. The homotopy

\begin{equation}
H_\tau(x,t)=\bigl(x,(1-\tau)t+\tau\rho(x)\bigr),
\qquad 0\leq\tau\leq1,
\label{mapsI:cw:eq:1.7}
\end{equation}

stays in \(T\). Indeed the set of heights at which a fixed \(x\) belongs to \(T\) is exactly \([j,\infty)\), where \(j\) is its first eligible stage, and both endpoint heights lie there. It fixes the section, starts at \(1_T\) and ends at \(s q\). This proves an explicit homotopy equivalence \(T\simeq U\).

No claim that the original manifold has been triangulated is needed. The maps \eqref{mapsI:cw:eq:1.6}--\eqref{mapsI:cw:eq:1.7} specify the exact comparison with a CW complex.

\subsection{Keep the original compact manifold and the actual punctured cobordism}
\label{mapsI:cw:cw-manifolds}

Let \(M\) be a compact smooth manifold without boundary. Choose finitely many charts \(x_i:U_i\to\mathbb R^d\), and smooth functions \(\phi_i\geq0\) whose supports lie inside \(U_i\) and whose positive sets cover \(M\). Finite smooth bump functions and a finite chart cover give this choice. The map

\begin{equation}
\iota:M\longrightarrow\mathbb R^{(d+1)k},\qquad
p\longmapsto\bigl(\phi_i(p),\,\phi_i(p)x_i(p)\bigr)_{i=1}^k
\label{mapsI:cw:eq:2.1}
\end{equation}

extends its blocks by zero off their charts and is smooth. It is injective: a positive block recovers both its chart and \(x_i(p)\) by division. Its derivative is injective: if it vanishes on a tangent vector, its first coordinate gives \(d\phi_i=0\), and then its second gives \(\phi_i\,dx_i=0\) for a positive block. Compactness and the Hausdorff target make the continuous inverse on the image continuous. Thus \eqref{mapsI:cw:eq:2.1} is a smooth embedding.

Use its Euclidean normal bundle, the orthogonal complement of \(d\iota(TM)\). The map

\begin{equation}
\mathcal L(p,v)=\iota(p)+v
\label{mapsI:cw:eq:2.2}
\end{equation}

has invertible derivative at every \((p,0)\): the tangent and normal summands map to their direct sum in the ambient Euclidean space. Here is the finite-dimensional inverse-function argument. After translations and an invertible linear change, write the map as \(F(x)=x+R(x)\) near zero, with \(\|DR\|<1/2\) on a closed ball. For \(y\) in a sufficiently small ball, \(x\mapsto y-R(x)\) maps the first ball into itself and is a contraction. Its iterates are uniformly Cauchy, giving the unique inverse value. The inequality \(\|F(x)-F(x^{\prime})\|\ge\|x-x^{\prime}\|/2\) proves continuity of the inverse. The remainder estimate in the differentiability formula then gives derivative \(DF(x)^{-1}\); induction gives smoothness. Thus the displayed map is a local diffeomorphism there. Compactness gives \(\epsilon>0\) on whose entire normal disk bundle it is a local diffeomorphism and injective. For injectivity, if no such radius existed, two distinct normal vectors with norms tending to zero would have equal images. Compactness gives convergent base subsequences; equality of their limits forces the same limiting base. The local inverse at that zero vector contradicts the equality for sufficiently late terms.

Its image \(U\) is open in Euclidean space. Write \(\mathcal L^{-1}(y)=(p(y),v(y))\). The maps

\begin{equation}
r(y)=p(y),\qquad
A_\tau(y)=\mathcal L\bigl(p(y),(1-\tau)v(y)\bigr)
\label{mapsI:cw:eq:2.3}
\end{equation}

give \(r\iota=1_M\) and a homotopy from \(1_U\) to \(\iota r\), fixed on \(\iota(M)\). All vectors in \eqref{mapsI:cw:eq:2.3} remain inside the chosen normal disk bundle.

Use the telescope and its maps from \cref{mapsI:cw:cw-open} for this actual \(U\). Define

\begin{equation}
\alpha=r q:T\longrightarrow M,\qquad
\beta=s\iota:M\longrightarrow T.
\label{mapsI:cw:eq:2.4}
\end{equation}

Then \(\alpha\beta=1_M\), and \(\beta\alpha=s\iota r q\) is homotopic first to \(s q\), through \(s A_\tau q\) with the parameter reversed, and then to \(1_T\) by reversing \eqref{mapsI:cw:eq:1.7}. This gives the CW homotopy type of the original \(M\), with both comparison maps exhibited.

For the original sphere construction take \(M=X\), with its original two disjoint coordinate disks. Write their given radii as \(R_0,R_1>0\), and retain

\begin{equation}
C=X\setminus\bigl(\operatorname{int}D_0\cup\operatorname{int}D_1\bigr).
\label{mapsI:cw:eq:2.5}
\end{equation}

There is also an explicit CW model of this actual compact manifold with boundary. Let \(W\subset X\) be the open set obtained by deleting the two closed concentric coordinate disks of radii \(R_i/2\). On the annulus in the \(i\)-th chart with radial coordinate \(u>R_i/2\), use

\begin{equation}
u\longmapsto(1-\tau)u+\tau\max(u,R_i),\qquad 0\leq\tau\leq1,
\label{mapsI:cw:eq:2.6}
\end{equation}

and keep the angular coordinate. Set the map equal to the identity outside these annuli. At \(u=R_i\) the formulas agree, and the annuli are disjoint. They give a continuous deformation retraction of \(W\) onto \(C\), fixed on \(C\); every intermediate radius stays greater than \(R_i/2\).

Restrict the normal disk bundle in \eqref{mapsI:cw:eq:2.2} to \(W\). Its image \(U_W\) is open in the ambient Euclidean space, and its normal contraction retracts it onto \(W\). Combining that contraction, \eqref{mapsI:cw:eq:2.6}, and the telescope construction on \(U_W\) gives explicit inverse homotopy comparisons between a CW complex and \(C\). The radii \(R_i\) in \eqref{mapsI:cw:eq:2.5} have not been changed: \(R_i/2\) defines an auxiliary open neighbourhood only.

\subsection{The cylinder extension and coherent simplex contractions}
\label{mapsI:cw:hurewicz-coherent}

For a closed unit disk \(D^m\) and \(I=[0,1]\), define

\begin{equation}
\lambda(v,t)=\min\left\{\frac{2}{2-t},\frac1{\lVert v\rVert}\right\},
\qquad
\mathcal R(v,t)=\bigl(\lambda(v,t)v,\,2+\lambda(v,t)(t-2)\bigr),
\label{mapsI:cw:eq:3.1}
\end{equation}

with the second argument of the minimum interpreted as \(+\infty\) at \(v=0\). The first choice of the minimum puts the result on the bottom; the second puts it on the side. Its time coordinate lies in \([0,t]\), and its radial coordinate has norm at most one. On the bottom and sides \(\lambda=1\), so \(\mathcal R\) fixes them. It is continuous, including at \(v=0\), because the bounded first term supplies the minimum throughout a sufficiently small neighbourhood of that point. Thus \eqref{mapsI:cw:eq:3.1} retracts the cylinder onto its bottom and sides.

A simplex is radially homeomorphic to a disk, from its barycentre, so this proves the same extension statement for simplex cylinders. In particular a map on the bottom and a compatible homotopy on the boundary extend to the cylinder by composition with \(\mathcal R\).

Let \(Y\) be path connected with basepoint \(y_0\), and suppose \(q\geq2\) and

\begin{equation}
\pi_i(Y,y_0)=0\qquad(1\leq i<q).
\label{mapsI:cw:eq:3.2}
\end{equation}

For every singular simplex of dimension at most \(q-1\), choose a homotopy to the constant simplex at \(y_0\), coherently on all faces. At vertices, choose paths to \(y_0\), with the constant path at \(y_0\). Suppose the lower-dimensional choices have been made. On the boundary of an \(i\)-simplex cylinder, prescribe its original simplex on the bottom, its constant map on the top, and the chosen face homotopies on the sides. They agree on intersections. This boundary is an \(i\)-sphere; for \(1\leq i<q\), \eqref{mapsI:cw:eq:3.2} makes its based class zero after transport to \(y_0\). A null-homotopy extends it across the cylinder. Choose that extension. For a simplex constant at \(y_0\), choose the constant extension.

For every \(q\)-simplex \(\sigma\), use \eqref{mapsI:cw:eq:3.1} to extend its bottom map and these chosen side homotopies. Its top \(\sigma'\) then has constant boundary. The quotient \(\Delta^q/\partial\Delta^q\), with the original orientation, is \(S^q\). Thus the top defines

\begin{equation}
a(\sigma)=[\sigma']\in\pi_q(Y,y_0).
\label{mapsI:cw:eq:3.3}
\end{equation}

The group is abelian for \(q\geq2\). Indeed concatenate cube representatives in either of two coordinate directions. The interchange law follows by dividing the square of those two coordinates into four rectangles. Both operations have the constant map as unit; applying interchange with two of the four entries equal to that unit proves that the operations agree and commute. Extend \eqref{mapsI:cw:eq:3.3} additively to singular \(q\)-chains.

We use the following finite sphere-pinch fact. If a triangulated \(q\)-sphere has its \((q-1)\)-skeleton mapped to the basepoint, its class is the sum of the classes of its oriented \(q\)-simplices. Here is a geometric proof. On each simplex collapse a collar of its boundary, and concentrate its map on a smaller interior disk; radial reparametrization gives a homotopy relative to the boundary. There are finitely many disjoint supporting disks. Remove the basepoint of the source sphere and use its Euclidean chart. For \(q\geq2\), the complement of finitely many disjoint small disks is path connected: a polygonal path can detour around each disk on its connected boundary. Move one sufficiently small support at a time along such a path to a prescribed sequence of disjoint slabs. A smooth vector field equal to the path's local translation near the moving support and cut off away from the other supports gives this isotopy on each short path segment; compactness divides the path into finitely many such segments. Radial shrinking and expanding supplies the required small radii. These isotopies preserve orientation. The final map is the concatenation of the individual disk maps, so its class is their sum, with the original simplex-orientation signs. The abelian law just proved removes any dependence on the chosen order.

\subsection{The first Hurewicz theorem, with inverse maps}
\label{mapsI:cw:hurewicz-theorem}

Define

\begin{equation}
h:\pi_q(Y,y_0)\longrightarrow H_q(Y;\mathbb Z),\qquad
h([f])=f_*[S^q].
\label{mapsI:cw:eq:4.1}
\end{equation}

The prism identity gives homotopy invariance. Under the two-disc pinch, the fundamental class of the sphere maps to the sum of the two positive disc generators. This follows from their local values and excision. Hence \(h\) is a homomorphism. It is natural under maps of based spaces.

\begin{theorem}
\label{mapsI:cw:hurewicz-theorem-result}
Under \eqref{mapsI:cw:eq:3.2}, \(h\) is an isomorphism and
\(\widetilde H_i(Y;\mathbb Z)=0\) for \(0\leq i<q\).
\end{theorem}

\begin{proof}
Apply the chosen simplex homotopies to the \(q\)-faces of a singular \((q+1)\)-simplex \(\tau\). They agree on their common faces. Formula \eqref{mapsI:cw:eq:3.1} extends their boundary homotopy across \(\tau\)'s cylinder. At its top, every \((q-1)\)-face is constant; its boundary sphere has class

\begin{equation}
\sum_{i=0}^{q+1}(-1)^i a(d_i\tau)
\label{mapsI:cw:eq:4.2}
\end{equation}

by the proved pinch fact. This sphere bounds the top \((q+1)\)-simplex, so \eqref{mapsI:cw:eq:4.2} is zero. Therefore \(a(\partial\tau)=0\), and the additive map \eqref{mapsI:cw:eq:3.3} induces

\begin{equation}
A:H_q(Y;\mathbb Z)\longrightarrow\pi_q(Y,y_0).
\label{mapsI:cw:eq:4.3}
\end{equation}

Triangulate each chosen simplex cylinder by the ordered prism triangulation. Its chain boundary is top minus bottom minus the corresponding face prisms. On a singular cycle \(z\), the face-prism terms cancel exactly: the homotopies of repeated singular faces are the same choices, and their coefficients are those of \(\partial z=0\). The top chain is consequently homologous to \(z\).

For the top maps work in the actual relative chain complex \(C_*(Y,\{y_0\};\mathbb Z)\). All constant simplices at \(y_0\) vanish in this quotient. The pair sequence identifies its positive-degree homology with \(H_*(Y;\mathbb Z)\): the point has no positive homology and its map on \(H_0\) is an isomorphism. Each top \(q\)-simplex has constant boundary and represents exactly \(h(a(\sigma))\) under this comparison, since the identity simplex is the positive disk-relative generator. Thus the prism argument proves

\begin{equation}
hA([z])=[z].
\label{mapsI:cw:eq:4.4}
\end{equation}

Conversely triangulate the source of a based sphere map \(f:S^q\to Y\), taking its basepoint to be a vertex and using a common vertex order on each face. The already chosen homotopies of its singular lower faces and \(q\)-simplices fit together into a homotopy of the entire sphere, fixing that basepoint. The final lower skeleton is constant. Its class is the oriented sum of the \(a\)-values by the same pinch fact. This is precisely \(A(h([f]))\), so

\begin{equation}
Ah([f])=[f].
\label{mapsI:cw:eq:4.5}
\end{equation}

Equations \eqref{mapsI:cw:eq:4.4}–\eqref{mapsI:cw:eq:4.5} prove the isomorphism. If \(i<q\), the coherent homotopies take an \(i\)-cycle entirely to constant simplices. The same prism cancellation makes its relative homology class zero. This proves positive-degree vanishing below \(q\); reduced \(H_0\) vanishes by path connectivity.
\end{proof}

For completeness, in degree one the Hurewicz map identifies \(H_1(Y;\mathbb Z)\) with the abelianization of \(\pi_1(Y,y_0)\). Choose paths \(p_y\) from \(y_0\) to each point \(y\). An edge \(\sigma\) from \(x\) to \(y\) determines \(p_x*\sigma*p_y^{-1}\). Sum these classes in the abelianized group over a cycle. The boundary of each singular triangle gives zero, since that triangle fills its edge loop and the inserted paths cancel. Changing the chosen paths contributes only endpoint terms, which cancel on a cycle. This defines the inverse: a based loop returns its original abelianized class, and the inserted paths cancel in homology. Both compositions are identities.

\subsection{The exact receiving sphere class and two exercises}
\label{mapsI:cw:cw-hurewicz-exercises}

Once \cref{mapsI:top:thm:h2-vanishing,mapsI:top:lem:cohomological-closure,mapsI:top:prop:geometric-pi1} have been applied, the original \(X\) satisfies
\(\pi_1(X)=0\), \(H_i(X;\mathbb Z)=0\) for \(1\leq i\leq5\), and
\(H_6(X;\mathbb Z)=\mathbb Z[X]\).
\cref{mapsI:cw:hurewicz-theorem-result} successively gives \(\pi_i(X)=0\) for \(2\leq i\leq5\), and then the actual isomorphism

\begin{equation}
h:\pi_6(X,x_0)\xrightarrow{\;\cong\;}H_6(X;\mathbb Z).
\label{mapsI:cw:eq:5.1}
\end{equation}

Choose a based representative \(g:S^6\to X\) of \(h^{-1}([X])\). Its degree is \(+1\), with
\(g_*[S^6]=[X]\). This constructs the original map used by the original sphere construction; no equivalence or smooth identification is presumed at this stage. The further homotopy-fibre comparison remains a separate receiving proof.

Equations \eqref{mapsI:cw:eq:1.6}--\eqref{mapsI:cw:eq:1.7} and \eqref{mapsI:cw:eq:2.3}--\eqref{mapsI:cw:eq:2.4} give all four maps and both inverse homotopies. The lower line is the separate integer Hurewicz isomorphism after the vanishing hypotheses have been established for the original \(X\), as in \eqref{mapsI:cw:eq:5.1}. It specifies the actual degree-one map without asserting smooth recognition.

\textbf{Exercise 5.1.}\label{mapsI:cw:exercise:5.1} In the telescope construction, prove that replacing \(\rho\) by another continuous function \(\rho'\) with \(x\in K_{\lfloor\rho'(x)\rfloor}\) changes the section by a homotopy, and give the homotopy.

\textbf{Solution.} If \(j(x)\) is the first eligible index, every allowed height belongs to \([j(x),\infty)\). Both \(\rho(x)\) and \(\rho'(x)\) lie there. The map
\begin{equation}
(x,\tau)\longmapsto
\bigl(x,(1-\tau)\rho(x)+\tau\rho'(x)\bigr)
\label{mapsI:cw:eq:5.2}
\end{equation}
therefore stays in the telescope. It is continuous in the subspace model \eqref{mapsI:cw:eq:1.4} because both functions are continuous. Projection by \(q\) remains the identity at every time. Thus the change is a homotopy through actual sections, with no relabelling of the cubical stages.

\textbf{Exercise 5.2.}\label{mapsI:cw:exercise:5.2} Replace the orientation class of \(X\) by \(-[X]\), keeping the original orientation of \(S^6\). Give the corresponding sphere map under \eqref{mapsI:cw:eq:5.1}, and compute its degree relative to each target orientation.

\textbf{Solution.} Let \(r:S^6\to S^6\) be the restriction of a linear reflection of \(\mathbb R^7\) whose fixed hyperplane contains the chosen sphere basepoint. In an orthonormal basis with that basepoint as the last vector, reflect the first coordinate. The map is therefore based, and its induced map on the outward-oriented sphere fundamental class is multiplication by \(-1\), by the determinant-sign and boundary-orientation calculation. Naturality of \eqref{mapsI:cw:eq:4.1} gives
\begin{equation}
h([g\circ r])=(g\circ r)_*[S^6]=-[X].
\label{mapsI:cw:eq:5.3}
\end{equation}
Thus \(g\circ r\) represents \(h^{-1}(-[X])\). It has degree \(-1\) relative to the original class \([X]\), and degree \(+1\) relative to the newly chosen class \(-[X]\). The map \(g\) and the original formula \(g_*[S^6]=[X]\) have not been silently changed.

\section{Path fibres and integral homotopy equivalences}
\label{mapsI:pf:section}

The CW and Hurewicz inputs are proved in \cref{mapsI:cw:section}; singular chains, excision, products and prisms are proved in \cref{mapsI:chain:section}. These are expanded classical proofs; the homological Whitehead conclusion is \cite[Corollary~4.33]{Hatcher2002}. All homology below is integral unless a coefficient group is displayed.

\subsection{Paths, lifting and the exact homotopy sequence}
\label{mapsI:pf:path-lifting}

For a space \(B\), give its path space \(B^I\), \(I=[0,1]\), the compact-open topology. Its subbasic sets are

\begin{equation}
[K,U]=\{\gamma:\gamma(K)\subset U\},
\label{mapsI:pf:eq:1.1}
\end{equation}

where \(K\subset I\) is compact and \(U\subset B\) is open. Evaluation is continuous: if \(\gamma(t)\in U\), choose a compact interval neighbourhood \(J\) of \(t\), relative to \(I\), with \(\gamma(J)\subset U\). The product \([J,U]\times\operatorname{int}_I J\) maps into \(U\). A map \(Z\to B^I\) is continuous exactly when its adjoint \(Z\times I\to B\) is continuous. For the reverse implication, a compact \(K\) is covered by finitely many product neighbourhoods on which the adjoint lies in \(U\); intersect their neighbourhoods in \(Z\). This proves continuity on \eqref{mapsI:pf:eq:1.1}, without a local-compactness assumption on \(Z\).

Given \(f:A\to B\), set

\begin{equation}
E_f=\{(a,\gamma)\in A\times B^I:\gamma(0)=f(a)\},\qquad
p_f(a,\gamma)=\gamma(1).
\label{mapsI:pf:eq:1.2}
\end{equation}

The inclusion \(i_f(a)=(a,c_{f(a)})\), where \(c_b\) is the constant path, and the projection \(r_f(a,\gamma)=a\) satisfy

\begin{equation}
r_fi_f=1_A,\qquad p_fi_f=f.
\label{mapsI:pf:eq:1.3}
\end{equation}

The homotopy \((a,\gamma)\mapsto(a,s\mapsto\gamma((1-u)s))\) starts at the identity and ends at \(i_fr_f\), fixing all constant paths.

We use a fibration to mean a map with the homotopy lifting property for every parameter space. The map \(p_f\) is such a fibration. Given \(H:Z\times I\to B\) and an initial lift \((a(z),\gamma_z)\), define

\begin{equation}
\eta_{z,t}(s)=
\begin{cases}
\gamma_z((1+t)s),&0\leq s\leq (1+t)^{-1},\\
H(z,(1+t)s-1),&(1+t)^{-1}\leq s\leq1.
\end{cases}
\label{mapsI:pf:eq:1.4}
\end{equation}

The two values agree at the join since \(\gamma_z(1)=H(z,0)\). The regions are closed and cover \(Z\times I\times I\); the pasting lemma and \eqref{mapsI:pf:eq:1.1} prove continuity. At time zero this is exactly the initial path, and its last endpoint is \(H(z,t)\). Thus \((a(z),\eta_{z,t})\) is the required lift.

Two relative lifting facts will be used. First, a lift can be prescribed on \(K\times\{0\}\cup L\times I\) for a CW pair \((K,L)\). For one cell this follows because

\begin{equation}
(D^j\times I,\ D^j\times\{0\}\cup S^{j-1}\times I)
\cong(D^j\times I,\ D^j\times\{0\}).
\label{mapsI:pf:eq:1.5}
\end{equation}

Here is the needed homeomorphism of pairs. The bottom-and-side cap is a disk: on \(v\in D^j\), put \(r=\|v\|\), \(\mu=\min(2,1/r)\), interpreting \(1/0=+\infty\), and send \(v\) to \((\mu v,2-\mu)\). For \(r\leq1/2\) this parametrizes the bottom disk; for \(r\geq1/2\) it parametrizes the side, with the rim mapped to the top rim. Match this rim parametrization with that of the top disk to obtain a boundary homeomorphism from the original cylinder boundary, taking its bottom to this cap. Extend it radially from an interior point of the convex cylinder. The case \(j=0\) only prescribes the initial point. The ordinary lifting property in these coordinates proves the cell step. Induct over cells; the weak topology gives a continuous lift. The same argument applies on products with \(I\): continuity can be tested by adjunction to the compact-open path space and then on every characteristic disk.

Second, this relative lifting also holds for the square with its bottom and two sides prescribed, with an arbitrary extra parameter space \(Z\). Reparametrize its boundary circle, piecewise linearly, taking the bottom edge onto the three-edge arc, and extend radially from its centre. Apply the ordinary lifting property with parameter \(Z\times I\). No CW hypothesis on \(Z\) is needed in this square case.

For a fibration \(p:E\to B\) and a path \(\alpha:b\leadsto b'\), lift \((e,t)\mapsto\alpha(t)\), starting at every \(e\in F_b=p^{-1}(b)\). Its endpoint is a map \(T_\alpha:F_b\to F_{b'}\). Two choices give homotopic endpoint maps: use the two lifts on the sides of the relative square, the initial inclusion on the bottom, and the constant-in-the-square-direction base path. A homotopy of paths relative their endpoints gives the same comparison. Concatenating lifts gives

\begin{equation}
T_{\alpha*\beta}\simeq T_\beta T_\alpha,\qquad
T_{c_b}\simeq1,\qquad
T_{\alpha^{-1}}T_\alpha\simeq1,\qquad
T_\alpha T_{\alpha^{-1}}\simeq1.
\label{mapsI:pf:eq:1.6}
\end{equation}

For the last two equations, contract a path followed by its reverse: traverse an initial segment and return along it. Path adjunction proves these homotopies continuous. Thus transport is a homotopy equivalence and gives the usual action of \(\pi_1(B,b)\) on \(H_*(F_b)\). In a simply connected base its action is the identity.

Pulling a fibration back along \(h:D\to B\) gives the fibration

\begin{equation}
h^*E=\{(d,e):h(d)=p(e)\}\longrightarrow D.
\label{mapsI:pf:eq:1.7}
\end{equation}

Lift a base homotopy through the second coordinate and retain its given first coordinate. If \(h_0\simeq h_1\), transport along that homotopy and its reverse gives inverse maps of their pullbacks up to homotopies over \(D\). Formula \eqref{mapsI:pf:eq:1.6}, applied with the \(D\)-coordinate held fixed, proves this assertion. A contraction of \(D^j\) therefore gives a fibre homotopy equivalence

\begin{equation}
(E|_{D^j},E|_{S^{j-1}})
\simeq(D^j\times F,S^{j-1}\times F).
\label{mapsI:pf:eq:1.8}
\end{equation}

It preserves every point of the base, so the same assertion holds for an annulus in place of the boundary sphere. The product-chain comparison then identifies

\begin{equation}
H_{j+q}(E|_{D^j},E|_{S^{j-1}};G)\cong H_q(F;G).
\label{mapsI:pf:eq:1.9}
\end{equation}

The oriented disk is the first factor. For completeness, its free relative chain complex contracts to \(\mathbb Z\) in degree \(j\): split its free cycle and boundary subgroups, using its sole relative homology generator. Tensoring those chain homotopies with the fibre complex and with \(G\) preserves them. This proves \eqref{mapsI:pf:eq:1.9} without assuming the fibre homology is free.

We next prove the exact sequence, including the disk comparison it uses. The relative group \(\pi_j(Y,D,d_*)\) consists of disk maps with boundary in \(D\) and a distinguished boundary point at \(d_*\), up to such homotopies. Equivalently use a cube with one distinguished face in \(D\) and all remaining faces at \(d_*\). Passing between the conventions contracts the remaining boundary cap while extending through a collar; the cylinder extension of \eqref{mapsI:pf:eq:1.5} provides that extension. Restriction to the distinguished face defines the connecting map. The sequence

\begin{equation}
\cdots\to\pi_j(D)\to\pi_j(Y)\to\pi_j(Y,D)
\xrightarrow{\partial}\pi_{j-1}(D)\to\pi_{j-1}(Y)\to\cdots
\label{mapsI:pf:eq:1.10}
\end{equation}

is exact. Here are its filling arguments. A relative disk whose boundary is null in \(D\) can have that null-homotopy attached along its boundary collar; capping that collar gives a based sphere in \(Y\) with the same relative class. A sphere in \(D\) dies in \(Y\) exactly when it bounds such a relative disk. Finally, if a sphere in \(Y\) is null as a relative disk, its relative null-homotopy compresses it into \(D\). To see the last assertion with its original boundary fixed, let \(F:D^j\times I\to Y\) be a relative null-homotopy. Its top and side lie in \(D\). Put

\begin{equation}
Q_s(v)=\bigl((1-s+s\mu)v,\ s(\mu-1)\bigr),
\qquad \mu=\min(2,1/\|v\|),\quad0\leq s\leq1.
\label{mapsI:pf:eq:1.11}
\end{equation}

This is the straight segment in the cylinder from \((v,0)\) to its top-and-side cap. It fixes every \((v,0)\) on the rim. Thus \(FQ_s\) compresses the original disk into \(D\) and fixes its full original boundary. These three fillings prove exactness at all group terms. The identical path fillings prove the pointed-set assertions at degrees one and zero. Disk concatenation proves that the displayed maps are homomorphisms where the terms are groups.

For \(F=p^{-1}(b_*)\), projection gives

\begin{equation}
p_*:\pi_j(E,F,e_*)\xrightarrow{\ \cong\ }\pi_j(B,b_*).
\label{mapsI:pf:eq:1.12}
\end{equation}

It is an isomorphism for \(j\geq2\) and a bijection of the relative pointed sets for \(j=1\). Indeed, in the cubical convention let \(J\) be all faces except the distinguished one. Deform the cube onto \(J\), fixing \(J\), then reverse that deformation and lift its image under a based sphere map into \(B\). Prescribe the constant lift on \(J\) by \eqref{mapsI:pf:eq:1.5}. The resulting disk is a lift representing \eqref{mapsI:pf:eq:1.12}. An explicit cap retraction, in \(D^{j-1}\times I\) coordinates when \(J\) is the top and sides, is

\begin{equation}
(v,t)\longmapsto\bigl(\lambda v,-1+\lambda(t+1)\bigr),
\qquad
\lambda=\min\left(\frac2{t+1},\frac1{\|v\|}\right).
\label{mapsI:pf:eq:1.13}
\end{equation}

Ignore the second entry at \(v=0\). Straight interpolation with the identity stays in the cylinder and fixes \(J\). If a projected relative disk is null, lift its based null-homotopy, keeping \(J\) constant. At its end the whole disk lies in \(F\), so \eqref{mapsI:pf:eq:1.11} makes its original relative class zero. This proves injectivity for groups. For two relative paths with homotopic projections, lift the homotopy square with the original paths on its sides and their distinguished endpoint on the third edge. Its fourth edge lies in \(F\), giving exactly a relative path homotopy. This proves injectivity of pointed sets.

Substitution in \eqref{mapsI:pf:eq:1.10} proves the fibration sequence

\begin{equation}
\cdots\to\pi_j(F)\to\pi_j(E)\xrightarrow{p_*}\pi_j(B)
\to\pi_{j-1}(F)\to\cdots\to\pi_0(E)\to\pi_0(B).
\label{mapsI:pf:eq:1.14}
\end{equation}

The connecting map takes the boundary of the lifted base disk. Projection and restriction show naturality directly. We use \eqref{mapsI:pf:eq:1.14} for \eqref{mapsI:pf:eq:1.2}, whose total space has the explicitly specified equivalence with \(A\) in \eqref{mapsI:pf:eq:1.3}.

\subsection{The filtered-chain construction and its convergence}
\label{mapsI:pf:filtered-chains}

Let \(C_*\) be a chain complex, with increasing subcomplexes \(F_pC_*\), zero for \(p<0\), whose union is \(C_*\). Write \(n=p+q\). For \(r\geq0\), set

\begin{equation}
\begin{aligned}
Z_r^{p,n}&=\{x\in F_pC_n:dx\in F_{p-r}C_{n-1}\},\\
B_r^{p,n}&=F_pC_n\cap d(F_{p+r}C_{n+1}).
\end{aligned}
\label{mapsI:pf:eq:2.1}
\end{equation}

In particular \(Z_0^{p,n}=F_pC_n\). For \(r\geq1\), define

\begin{equation}
E_r^{p,q}=
\frac{Z_r^{p,p+q}}
{Z_{r-1}^{p-1,p+q}+B_{r-1}^{p,p+q}},
\qquad
d_r[x]=[dx]\in E_r^{p-r,q+r-1}.
\label{mapsI:pf:eq:2.2}
\end{equation}

The denominator lies in the numerator. For its first term, \(dz\in F_{p-r}\); its second term consists of cycles. A boundary \(dx\) lies in the target numerator because its next boundary is zero. Replacing \(x\) by the first denominator term changes \(dx\) by a target \(B_{r-1}\), since that term lies in \(F_{p-1}=F_{(p-r)+(r-1)}\). Replacing it by a boundary has no effect. Thus \(d_r\) is defined and \(d_r^2=0\).

We verify the page transition, rather than assume it. If \(d_r[x]=0\), write

\begin{equation}
dx=z+dy,\quad
z\in Z_{r-1}^{p-r-1,n-1},\quad
y\in F_{p-1}C_n,\quad dy\in F_{p-r}C_{n-1}.
\label{mapsI:pf:eq:2.3}
\end{equation}

Then \(y\in Z_{r-1}^{p-1,n}\), so \(x-y\) has the same page-\(r\) class and belongs to \(Z_{r+1}^{p,n}\). Conversely every such representative is a page-\(r\) cycle. Images of \(d_r\) are represented by \(B_r^{p,n}\), since a chain in \(F_{p+r}\) whose boundary lies in \(F_p\) is exactly a source numerator. In the kernel, the remaining ambiguity in \(F_{p-1}\) has boundary in \(F_{p-r-1}\); it is \(Z_r^{p-1,n}\). Consequently

\begin{equation}
H(E_r,d_r)^{p,q}
=Z_{r+1}^{p,n}/(Z_r^{p-1,n}+B_r^{p,n})
=E_{r+1}^{p,q}.
\label{mapsI:pf:eq:2.4}
\end{equation}

At the first page \eqref{mapsI:pf:eq:2.2} is precisely \(H_{p+q}(F_pC/F_{p-1}C)\).

Assume these first-page groups vanish for \(q<0\). The subsequent pages do too, by \eqref{mapsI:pf:eq:2.4}. For a fixed \((p,q)\) in the first quadrant, an outgoing differential is zero if \(r>p\), and an incoming one is zero if \(r>q+1\). Thus the entry stabilizes after \(r>\max(p,q+1)\). For \(r>p\), the numerator in \eqref{mapsI:pf:eq:2.2} is the actual cycle subgroup in \(F_pC_n\), and its first denominator consists of the actual cycles in \(F_{p-1}C_n\). Exhaustiveness gives

\begin{equation}
\bigcup_r B_r^{p,n}=F_pC_n\cap dC_{n+1}.
\label{mapsI:pf:eq:2.5}
\end{equation}

Passing to that union in a stabilized entry proves

\begin{equation}
E_\infty^{p,n-p}\cong F_pH_n(C)/F_{p-1}H_n(C),
\qquad
F_pH_n(C)=\operatorname{im}\bigl(H_n(F_pC)\to H_n(C)\bigr).
\label{mapsI:pf:eq:2.6}
\end{equation}

This filtration is exhaustive, since a cycle is a finite chain in some \(F_pC\). For \(p>n\) its consecutive quotient is zero. Hence exhaustiveness also gives

\begin{equation}
0=F_{-1}H_n\subset F_0H_n\subset\cdots\subset F_nH_n=H_n.
\label{mapsI:pf:eq:2.7}
\end{equation}

No hypothesis \(C_n\subset F_nC_n\) was used. Equations \eqref{mapsI:pf:eq:2.6}--\eqref{mapsI:pf:eq:2.7} mean convergence with a finite filtration; they do not identify the group with a direct sum of its quotients. A filtration-preserving chain map preserves \eqref{mapsI:pf:eq:2.1}, \eqref{mapsI:pf:eq:2.2} and \eqref{mapsI:pf:eq:2.6}, proving all the naturality used below.

\subsection{The fibration filtration and both actual edges}
\label{mapsI:pf:serre-homology}

Let \(p:E\to B\) be a fibration, with \(B\) a path-connected, simply connected CW complex and \(F=p^{-1}(b_*)\) path connected; choose \(b_*\) to be a vertex. Orient the cells of \(B\). For an abelian group \(G\), we will prove

\begin{equation}
E_2^{p,q}=H_p(B;H_q(F;G)),\qquad
d_r:E_r^{p,q}\longrightarrow E_r^{p-r,q+r-1}
\quad\Longrightarrow\quad H_{p+q}(E;G).
\label{mapsI:pf:eq:3.1}
\end{equation}

The coefficient system is constant because of \eqref{mapsI:pf:eq:1.6} and simple connectivity. Convergence has exactly the meaning \eqref{mapsI:pf:eq:2.6}--\eqref{mapsI:pf:eq:2.7}.

First, a compact subset of a CW complex is contained in a finite subcomplex. Otherwise choose one point in each of infinitely many distinct cells met by the subset. The closure of any cell meets only finitely many cells, by closure finiteness. Its intersection with any subset of the chosen points is therefore finite and closed. The weak topology says that every subset of those points is closed in \(B\). They form an infinite closed discrete subspace of the original compact subset, which is impossible. The finitely many cells that remain, together with their finite collections of boundary cells, give the asserted finite subcomplex.

Put \(E_p=p^{-1}(B^p)\) and

\begin{equation}
F_pC_*(E;G)=C_*(E_p;G),\qquad E_{-1}=\varnothing.
\label{mapsI:pf:eq:3.2}
\end{equation}

The projected image of a singular simplex is compact, so the preceding argument puts it in some finite skeleton. Thus \eqref{mapsI:pf:eq:3.2} is exhaustive. The first page of \cref{mapsI:pf:filtered-chains} is \(H_{p+q}(E_p,E_{p-1};G)\).

We calculate these relative groups with the actual attaching disks. In \(B^p\), let \(N\) be \(B^{p-1}\) together with the regions \(\|v\|>1/3\) of its characteristic \(p\)-disks, and let \(N_0\) use \(\|v\|>2/3\). They are open in \(B^p\) by the weak topology. The skeleton together with the regions \(\|v\|\geq2/3\) is closed, contains \(N_0\), and lies in \(N\), so \(\overline{N_0}\subset N\). Expanding each annulus radially to its outer sphere gives a deformation of \(N\) onto \(B^{p-1}\), fixed on that skeleton. It is continuous on characteristic disks and on their products with \(I\). Lift it, starting with the identity on \(p^{-1}(N)\). Its last map lands in \(E_{p-1}\); on \(E_{p-1}\) the whole lift stays in \(E_{p-1}\). Thus \(E_{p-1}\hookrightarrow p^{-1}(N)\) is a homotopy equivalence, even though the lifted homotopy need not fix its points. The pair sequence gives

\begin{equation}
H_*(E_p,E_{p-1};G)\cong H_*(E_p,p^{-1}(N);G).
\label{mapsI:pf:eq:3.3}
\end{equation}

Excision now removes \(p^{-1}(N_0)\), whose closure lies in \(p^{-1}(N)\). What remains is the disjoint union, over the open \(p\)-cells, of the disk cores \(\|v\|\leq2/3\) paired with \(1/3<\|v\|\leq2/3\). This is a topological disjoint union: a union of closed subsets of these cores is closed in \(B^p\), as can be tested on every characteristic disk. Its pullback has the same disjoint-union description. The characteristic maps are homeomorphisms on these cores. Trivialize each pullback up to fibre homotopy by \eqref{mapsI:pf:eq:1.8}, then replace the annulus by its outer sphere, using its radial retraction and the pair sequence. Formula \eqref{mapsI:pf:eq:1.9} gives

\begin{equation}
H_{p+q}(E_p,E_{p-1};G)
\cong\bigoplus_{\text{\(p\)-cells }e}H_q(F;G).
\label{mapsI:pf:eq:3.4}
\end{equation}

Transport to the reference fibre fixes the coefficient identification. At \(p=0\) the formula is the homology of the disjoint union of the vertex fibres. Every chain has finite support, which accounts for the direct sum, including when there are infinitely many cells. Negative \(q\) gives zero. \cref{mapsI:pf:filtered-chains} therefore supplies a first-quadrant sequence with finite-filtration convergence.

We identify its first differential completely. It is the pair boundary followed by the next relative quotient:

\begin{equation}
H_{p+q}(E_p,E_{p-1};G)\xrightarrow{\partial}
H_{p+q-1}(E_{p-1};G)\longrightarrow
H_{p+q-1}(E_{p-1},E_{p-2};G).
\label{mapsI:pf:eq:3.5}
\end{equation}

Over one characteristic disk, a class is its positive relative disk generator \(c\) crossed with a fibre cycle \(z\). The full boundary formula is

\begin{equation}
\partial(c\times z)=\partial c\times z+(-1)^p c\times\partial z
=\partial c\times z.
\label{mapsI:pf:eq:3.6}
\end{equation}

Consequently \eqref{mapsI:pf:eq:3.5} is the lifted attaching sphere with the same fibre coefficient.

Here is the local degree calculation of each of its components. For \(p\geq2\), put \(m=p-1\) and fix a target open \(m\)-cell. Choose concentric balls \(V_0\Subset V_1\Subset V_2\) inside its coordinate chart. Triangulate the compact attaching \(S^m\) finely enough that stars meeting the inverse image of \(\overline{V_2}\) are mapped into that chart after enlarging \(V_2\) slightly. On these stars, interpolate the coordinate values at vertices to obtain a piecewise affine map uniformly close to the original map. Choose a piecewise linear cutoff equal to one on a neighbourhood of the inverse image of \(\overline{V_1}\), supported in the inverse image of the enlarged chart region. Such a cutoff is obtained by assigning zero and one on vertices after a further subdivision, with the intervening stars lying between the two closed sets. Replace the coordinate map there by its interpolation multiplied by this cutoff plus the original map multiplied by one minus the cutoff. The error can be made smaller than each separation between the chosen balls. Straight interpolation gives a homotopy in the same chart, fixed outside it. Every inverse image of \(V_0\) then lies in the region where the modified map is piecewise affine.

Choose a point of \(V_0\) avoiding the images of all lower-dimensional faces and all deficient-rank affine pieces. There are finitely many of these proper affine subsets. Its inverse images are finitely many points in invertible affine pieces. Outside small neighbourhoods of those points the compact source has image missing the chosen point. A sufficiently small target ball therefore has inverse image precisely a disjoint union of source disks, each mapped homeomorphically onto that ball, with sign \(+1\) or \(-1\). The sum of the signs is the cellular incidence coefficient: excision maps the attaching sphere to the target relative disk group and counts its oriented generators. This also proves this local description of degree, rather than requiring a smooth transversality assertion.

Lift the homotopy of the attaching map with the fibre as additional parameter. It does not change \eqref{mapsI:pf:eq:3.5}, because it is a homotopy within \(E_{p-1}\). Over each of the resulting source disks, use the two fibre homotopy trivializations. The map takes the form

\begin{equation}
(v,z)\longmapsto(\varphi(v),\psi_v(z)).
\label{mapsI:pf:eq:3.7}
\end{equation}

Contract \(v\) in the second coordinate of \eqref{mapsI:pf:eq:3.7}, retaining its first coordinate. This is a homotopy of pairs: its relative boundary condition depends on the first coordinate alone. Its effect on relative homology is thus the local sign of \(\varphi\) multiplied by the homology map of \(\psi_{v_0}\). That map is the transport comparison on the two fibres. To verify it, apply the map of pullback fibrations to a chosen lifted path and compare it with the other chosen lift by the relative square in \cref{mapsI:pf:path-lifting}. Under transport to the reference fibre it is the identity. A different choice inserts a loop of \(B\), which is null-homotopic, so \eqref{mapsI:pf:eq:1.6} removes that ambiguity. Summing the local signs proves that \eqref{mapsI:pf:eq:3.5} is the cellular incidence coefficient times the original fibre class.

For \(p=1\), the boundary of the oriented interval is its terminal point minus its initial point. Transport identifies the two coefficient maps with the identity and gives precisely those two signed vertex terms. For \(p=0\) the differential is zero. Hence

\begin{equation}
(E_1^{*,q},d_1)=
(C_*^{\mathrm{cell}}(B;H_q(F;G)),\partial_{\mathrm{cell}}).
\label{mapsI:pf:eq:3.8}
\end{equation}

The cellular complex with any constant coefficient group \(A\) computes singular homology. This follows here from the same calculation with the identity fibration of \(B\): the relative disk homology is concentrated in row zero, with coefficient \(A\); \eqref{mapsI:pf:eq:3.5} is exactly its cellular boundary; the filtered-chain sequence has no possible differential after its second page. Its sole quotient in total degree \(n\) is therefore \(H_n(B;A)\), by \eqref{mapsI:pf:eq:2.6}--\eqref{mapsI:pf:eq:2.7}. This proves the assertion without an additional cellular comparison theorem. Taking \(A=H_q(F;G)\) proves \eqref{mapsI:pf:eq:3.1}.

We need the actual homology maps at both edges. First,

\begin{equation}
E_\infty^{0,n}=F_0H_n(E;G)
=\operatorname{im}\bigl(H_n(F;G)\to H_n(E;G)\bigr).
\label{mapsI:pf:eq:3.9}
\end{equation}

Indeed \(E_0\) is the union of vertex fibres. Transport along a path from any vertex to \(b_*\), followed by inclusion into \(E\), is homotopic in \(E\) to that vertex fibre's inclusion, by its defining lifted path. Thus all vertex fibres have the same image. On page two the entry is \(H_n(F;G)\); later pages take its quotients by incoming images, since no differential leaves column zero.

For projection, use the map of filtered complexes induced by \(p:E\to B\), with both filtrations taken over the same CW base. It maps the given fibration to the identity fibration. On row zero of the second page it is the identity of \(H_n(B;G)\), since \(H_0(F;G)=G\). No differential enters that row for \(r\geq2\), so every subsequent bottom-row entry is a subgroup of its predecessor. Naturality from \cref{mapsI:pf:filtered-chains} identifies its final subgroup with the image of projection. More explicitly, projection kills \(F_{n-1}H_n(E;G)\), since its image factors through \(H_n(B^{n-1};G)=0\), as follows from the just-proved cellular calculation. On the final quotient it is the inclusion \(E_\infty^{n,0}\hookrightarrow H_n(B;G)\). We have proved

\begin{equation}
\begin{aligned}
\operatorname{im}\bigl(p_*:H_n(E;G)\to H_n(B;G)\bigr)
&=E_\infty^{n,0}\subset E_2^{n,0},\\
\ker p_*&=F_{n-1}H_n(E;G).
\end{aligned}
\label{mapsI:pf:eq:3.10}
\end{equation}

These are the two edges needed below. They include the possible incoming transgression to the fibre term; no assertion of automatic survival has been made.

\subsection{From weak equivalences to inverse maps}
\label{mapsI:pf:cw-inverses}

A weak equivalence is a map inducing a bijection on path components and isomorphisms on all based positive homotopy groups. We prove the inverse-map assertion for connected CW complexes without changing the original map to a cellular one.

First suppose \((Y,D)\) has \(D\) path connected, every component of \(Y\) meeting \(D\), and all relative homotopy groups zero. Every map \(u:(K,L)\to(Y,D)\) from a CW pair can be homotoped, relative \(L\), into \(D\). Choose paths into \(D\) for the vertices outside \(L\). Extend their homotopy over all higher cells using the cylinder retraction of \cref{mapsI:cw:cw-manifolds}. Suppose the preceding relative skeleton already maps into \(D\). A characteristic \(j\)-disk has boundary in \(D\); its zero relative class supplies a relative null-homotopy, and \eqref{mapsI:pf:eq:1.11} turns it into a homotopy into \(D\) fixing its full boundary. For \(j=1\), the same relative path statement applies. Perform this on all \(j\)-cells and extend it to the rest of \(K\), again by cylinder extension.

Use the time interval \([1-2^{-j},1-2^{-(j+1)}]\) for the \(j\)-dimensional step, starting with \(j=0\). Later steps fix each earlier relative skeleton, and fix \(L\) throughout. On a fixed characteristic \(d\)-disk the homotopy is stationary after the \(d\)-dimensional step. Its restriction to that disk times \(I\), including time one, is therefore continuous. To verify continuity of the whole homotopy, take its adjoint \(K\to Y^I\). The restrictions to all characteristic disks are continuous by the compact-open adjunction; the weak topology makes the adjoint continuous on \(K\), and evaluation makes the original homotopy continuous. Its time-one map lies in \(D\). This proves the assertion also for infinite-dimensional \(K\).

Now let \(f:A\to B\) be a weak equivalence of path-connected spaces. Form its actual mapping cylinder

\begin{equation}
M_f=(A\times I)\amalg B\,/\bigl((a,1)\sim f(a)\bigr),
\qquad j(a)=[a,0].
\label{mapsI:pf:eq:4.1}
\end{equation}

The retraction \(r:M_f\to B\) sends \([a,t]\) to \(f(a)\) and fixes \(B\). Moving \(t\) to \(t+u(1-t)\) gives a deformation retraction onto \(B\). Thus \(rj=f\) makes \(j\) a weak equivalence. Exactness of \eqref{mapsI:pf:eq:1.10} gives \(\pi_k(M_f,j(A))=0\) at every positive degree, with the path statement at degree one as well. The preceding compression consequently shows, for every CW complex \(K\), that

\begin{equation}
f_*:[K,A]\xrightarrow{\ \cong\ }[K,B],
\label{mapsI:pf:eq:4.2}
\end{equation}

where brackets mean unbased homotopy classes. Surjectivity follows by viewing a map \(K\to B\) in \(M_f\), compressing it into \(j(A)\), and retracting the resulting homotopy by \(r\). For injectivity, a homotopy between \(f a_0\) and \(f a_1\) in \(B\), together with the two cylinder paths in \eqref{mapsI:pf:eq:4.1}, gives a homotopy between \(j a_0\) and \(j a_1\) in \(M_f\). Apply compression to the CW pair \((K\times I,K\times\partial I)\), fixing both endpoint maps. The compressed homotopy lies in \(j(A)\) and proves \(a_0\simeq a_1\). The product is a CW pair because \(I\) has its finite interval cell structure; its characteristic prisms give the product weak topology by the same adjunction argument.

When \(A\) and \(B\) are CW complexes, apply surjectivity of \eqref{mapsI:pf:eq:4.2} with \(K=B\) to obtain \(h:B\to A\) with \(fh\simeq1_B\). Then \(f(hf)\simeq f\), so injectivity with \(K=A\) gives \(hf\simeq1_A\). This constructs a homotopy inverse to the given \(f\).

If the original spaces have specified CW comparisons

\begin{equation}
\alpha_A:T_A\to A,\quad\beta_A:A\to T_A,\qquad
\alpha_B:T_B\to B,\quad\beta_B:B\to T_B,
\label{mapsI:pf:eq:4.3}
\end{equation}

with their specified inverse homotopies, apply the construction to

\begin{equation}
\bar f=\beta_B f\alpha_A:T_A\to T_B.
\label{mapsI:pf:eq:4.4}
\end{equation}

If \(\bar h\) is its inverse, the actual inverse on the original spaces is

\begin{equation}
h=\alpha_A\bar h\beta_B:B\to A.
\label{mapsI:pf:eq:4.5}
\end{equation}

Indeed \(\alpha_B\bar f\simeq f\alpha_A\) and \(\bar f\beta_A\simeq\beta_Bf\), by the two homotopies in \eqref{mapsI:pf:eq:4.3}. Substituting them in \(fh\) and \(hf\), followed by \(\bar f\bar h\simeq1\) and \(\bar h\bar f\simeq1\), gives \(fh\simeq1_B\) and \(hf\simeq1_A\). This proves the required correspondence, with both domains and codomains retained.

\subsection{The integer comparison and the original sphere maps}
\label{mapsI:pf:integral-comparison}

\begin{theorem}
\label{mapsI:pf:theorem:5.1}
Let \(f:A\to B\) be an integral homology equivalence between path-connected, simply connected spaces with specified CW homotopy comparisons. Then \(f\) is a homotopy equivalence. The construction below gives the inverse through those comparisons, rather than replacing \(f\) by an unspecified map.
\end{theorem}

\begin{proof}
Apply \eqref{mapsI:pf:eq:4.4}. The resulting map between CW complexes is again an integral homology equivalence and both CW complexes are simply connected. It is enough to prove that this map is a weak equivalence, since \cref{mapsI:pf:cw-inverses} then supplies \eqref{mapsI:pf:eq:4.5}. Denote the two CW complexes for this proof by \(A'\) and \(B'\), their map by \(f'\), its path-fibration total space by \(E\), and its fibre at the chosen vertex by \(K\). The retraction \eqref{mapsI:pf:eq:1.3} makes \(E\) homotopy equivalent to \(A'\), and \(p_*:H_j(E)\to H_j(B')\) is an isomorphism for every \(j\).

The fibre \(K\) is path connected. To check the component assertion directly, connect any two fibre points by a path in the connected total space \(E\). Its projection is a loop in the simply connected \(B'\), so it contracts relative its endpoints. Lift this homotopy with the original path on its bottom and the two constant endpoint lifts on its sides, using the relative square. Its top is a path in \(K\) connecting the original fibre points.

Choose \(a_*\in A'\) and a path from \(f'(a_*)\) to the reference vertex of \(B'\); together they specify a point \(e_*\in K\). The retraction \(r_f\) sends that point to \(a_*\). The basepoint comparison from the constant path to the chosen path projects to exactly this chosen base path. The degree-two Hurewicz maps for \(A'\) and \(B'\) are isomorphisms and are natural, by \cref{mapsI:cw:hurewicz-theorem-result}. The homology isomorphism \(f'_*\), with this basepoint transport, consequently makes \(\pi_2(E,e_*)\to\pi_2(B',b_*)\) an isomorphism. In \eqref{mapsI:pf:eq:1.14}, its surjectivity and \(\pi_1(E)=0\) imply

\begin{equation}
\pi_1(K)=0.
\label{mapsI:pf:eq:5.1}
\end{equation}

Suppose \(K\) has a nonzero positive homotopy group, and let \(n\geq2\) be its least degree. First Hurewicz, which was proved for the actual path-connected space rather than only for CW spaces, gives

\begin{equation}
H_i(K)=0\quad(0<i<n),\qquad
H_n(K)\cong\pi_n(K)\ne0.
\label{mapsI:pf:eq:5.2}
\end{equation}

Apply \eqref{mapsI:pf:eq:3.1} to \(E\to B'\). For the entry \((0,n)\), an incoming page-\(r\) source is \((r,n-r+1)\). It is zero for \(2\leq r\leq n\), by \eqref{mapsI:pf:eq:5.2}. For \(r>n+1\) its second coordinate is negative. Thus the only possible incoming differential is

\begin{equation}
d_{n+1}:E_{n+1}^{n+1,0}\longrightarrow E_{n+1}^{0,n}.
\label{mapsI:pf:eq:5.3}
\end{equation}

There is no outgoing differential from column zero. The source in \eqref{mapsI:pf:eq:5.3} is exactly \(H_{n+1}(B')\) at that page: bottom-row entries have no incoming maps, and each earlier outgoing target from \((n+1,0)\) has fibre degree \(r-1\) between one and \(n-1\), hence is zero. For later pages its target column is negative. Therefore

\begin{equation}
E_\infty^{n+1,0}=\ker d_{n+1}\subset H_{n+1}(B').
\label{mapsI:pf:eq:5.4}
\end{equation}

By \eqref{mapsI:pf:eq:3.10} this subgroup is the image of \(p_*:H_{n+1}(E)\to H_{n+1}(B')\). That map is surjective. Thus the kernel in \eqref{mapsI:pf:eq:5.4} is the entire source, and \(d_{n+1}=0\). It follows that

\begin{equation}
H_n(K)=E_\infty^{0,n}=F_0H_n(E).
\label{mapsI:pf:eq:5.5}
\end{equation}

Since \(n\geq2\), this subgroup lies in \(F_{n-1}H_n(E)=\ker p_*\), by \eqref{mapsI:pf:eq:3.10}. But \(p_*:H_n(E)\to H_n(B')\) is injective. Thus \eqref{mapsI:pf:eq:5.5} is zero, contradicting \eqref{mapsI:pf:eq:5.2}.

There is consequently no first nonzero positive homotopy group: all of them vanish. Together with path connectivity, \eqref{mapsI:pf:eq:1.14} says that \(p\), and hence \(f'\), induces every homotopy isomorphism. \cref{mapsI:pf:cw-inverses} gives its inverse \(\bar h\) and the original inverse \eqref{mapsI:pf:eq:4.5}. Both homotopies there have already been proved. This completes the theorem.
\end{proof}

\emph{\cref{mapsI:pf:theorem:5.1}, equations \eqref{mapsI:pf:eq:5.2}--\eqref{mapsI:pf:eq:5.5}. The entry \(H_n(K)\) may receive \(d_{n+1}\). Surjectivity in degree \(n+1\) first makes that map zero. Injectivity in degree \(n\) then makes the surviving fibre subgroup zero. These are the actual proved homology edge maps.}

We now apply the theorem with the exact maps and orientation classes of \cref{mapsI:course:homotopy-section}. The already proved integral and fundamental-group calculations give \(\pi_1(X)=0\) and

\begin{equation}
H_k(X;\mathbb Z)=
\begin{cases}\mathbb Z,&k=0,6,\\0,&k\ne0,6,\end{cases}
\qquad [X]\in H_6(X;\mathbb Z)
\label{mapsI:pf:eq:5.6}
\end{equation}

with the original complex orientation. Successive applications of first Hurewicz give \(\pi_2(X)=\cdots=\pi_5(X)=0\), then an isomorphism \(h_6:\pi_6(X)\to H_6(X;\mathbb Z)\). Choose \(g:S^6\to X\) representing \(h_6^{-1}[X]\). Its positive sphere class satisfies

\begin{equation}
g_*[S^6]=[X].
\label{mapsI:pf:eq:5.7}
\end{equation}

This makes \(g\) an integral homology equivalence in every degree. Retain the actual maps of \cref{mapsI:cw:cw-manifolds},

\begin{equation}
\alpha_X:T_X\to X,\quad\beta_X:X\to T_X,\qquad
\alpha_X\beta_X=1_X,\quad\beta_X\alpha_X\simeq1_{T_X}.
\label{mapsI:pf:eq:5.8}
\end{equation}

For \(\bar g=\beta_Xg:S^6\to T_X\), \cref{mapsI:pf:theorem:5.1} and \cref{mapsI:pf:cw-inverses} supply \(\bar h:T_X\to S^6\). The inverse to the original \(g\) is

\begin{equation}
h_X=\bar h\beta_X:X\to S^6,\qquad
h_Xg\simeq1_{S^6},\qquad gh_X\simeq1_X.
\label{mapsI:pf:eq:5.9}
\end{equation}

For the second equality, use \(g=\alpha_X\bar g\) exactly, then \(\bar g\bar h\simeq1_{T_X}\), then \(\alpha_X\beta_X=1_X\). Thus the orientation degree in \eqref{mapsI:pf:eq:5.7} has not been replaced or rescaled.

For the cobordism, retain the two original disjoint orientation-preserving coordinate disks, their radii \(R_0,R_1\), and

\begin{equation}
C=X\setminus(\operatorname{int}D_0\cup\operatorname{int}D_1),
\qquad i_j:\partial D_j\hookrightarrow C.
\label{mapsI:pf:eq:5.10}
\end{equation}

Excision for the original two discs and the integral homology of \(X\) give the actual pair sequence

\begin{equation}
0\longrightarrow H_6(C)\longrightarrow
\mathbb Z[X]\xrightarrow{\,1\mapsto(1,1)\,}
\mathbb Z[D_0,\partial D_0]\oplus\mathbb Z[D_1,\partial D_1]
\xrightarrow{\partial}H_5(C)\longrightarrow0.
\label{mapsI:pf:eq:5.11}
\end{equation}

Consequently \(H_6(C)=0\), \(H_5(C)=\mathbb Z^2/\mathbb Z(1,1)\), and the two actual boundary classes are the opposite generators \([(1,0)]\) and \([(0,1)]=-[(1,0)]\). The remaining positive homology groups vanish. Van Kampen gives \(\pi_1(C)=0\), by adding the two simply connected discs back along collars of their simply connected boundary spheres, and each boundary is a simply connected \(S^5\). Each \(i_j\) is therefore an integral homology equivalence, with its own displayed sign. The CW model of \cref{mapsI:cw:cw-manifolds} for \(C\) preserves the radii in \eqref{mapsI:pf:eq:5.10}; write its maps as \(\alpha_C,\beta_C\). Applied to \(\beta_Ci_j\), the same construction yields

\begin{equation}
h_j=\bar h_j\beta_C:C\to\partial D_j,\qquad
h_ji_j\simeq1_{\partial D_j},\qquad i_jh_j\simeq1_C.
\label{mapsI:pf:eq:5.12}
\end{equation}

This proves that the particular punctured manifold in \eqref{mapsI:pf:eq:5.10}, with its original two inclusions, is an \(h\)-cobordism. Equations \eqref{mapsI:pf:eq:5.9} and \eqref{mapsI:pf:eq:5.12} complete the homotopy assertions. Its smooth product structure and the extension problem for the later gluing diffeomorphism require the smooth recognition arguments; the homotopy comparison alone does not supply them.

\subsection{Two worked checks of the mechanism}
\label{mapsI:pf:path-fibre-exercises}

\textbf{Exercise 6.1 --- The incoming map cannot be omitted.}\label{mapsI:pf:exercise:6.1} Fix \(n\geq2\) and an integer \(m\). Take the free chain complex

\begin{equation}
C_0=\mathbb Ze,\quad C_n=\mathbb Zb,\quad C_{n+1}=\mathbb Za,
\qquad da=mb,\quad db=de=0.
\label{mapsI:pf:eq:6.1}
\end{equation}

Put \(e,b\) in filtration zero, \(a\) in filtration \(n+1\), and use the increasing filtration they generate. Compute its pages and both relevant homology maps under the quotient \(P:C\to C/\mathbb Zb\). This is an algebraic model of the edge calculation; no fibration realization is assumed.

\textbf{Solution.} The first and second pages have \(\mathbb Z\) at \((0,0),(0,n),(n+1,0)\) and zero elsewhere. The boundary \(da\) drops filtration by \(n+1\), so every differential before page \(n+1\) is zero. Formula \eqref{mapsI:pf:eq:2.2} then gives exactly

\begin{equation}
d_{n+1}:\mathbb Z[a]\longrightarrow\mathbb Z[b],
\qquad [a]\longmapsto m[b].
\label{mapsI:pf:eq:6.2}
\end{equation}

The final entries are \(\ker(m:\mathbb Z\to\mathbb Z)\) at \((n+1,0)\), \(\mathbb Z/m\mathbb Z\) at \((0,n)\), and \(\mathbb Z\) at \((0,0)\). These agree directly with the homology of \eqref{mapsI:pf:eq:6.1}.

The quotient has \(H_{n+1}(C/\mathbb Zb)=\mathbb Z[a]\) and \(H_n(C/\mathbb Zb)=0\). Its degree-\(n+1\) map has image \(\ker m\), precisely the bottom-row subgroup. If \(m\ne0\), that image is zero and is not surjective. If \(m=0\), it is surjective and the incoming differential is zero; but the degree-\(n\) map has kernel \(\mathbb Z[b]\) and is not injective. For \(m=1\) or \(m=-1\), the nonzero fibre-row entry is killed completely by the incoming differential. For all other nonzero \(m\), its surviving group is exactly \(\mathbb Z/m\mathbb Z\), retaining the full integer \(m\) and its sign in \eqref{mapsI:pf:eq:6.2}. This verifies why both adjacent-degree arguments in \cref{mapsI:pf:theorem:5.1} are needed.

\textbf{Exercise 6.2 --- The looped five- and six-spheres.}\label{mapsI:pf:exercise:6.2} Determine the additive integral homology of \(\Omega S^d\), for \(d=5\) and \(d=6\), using the proved filtration.

\textbf{Solution.} Use \(f:\{b_*\}\to S^d\) in \eqref{mapsI:pf:eq:1.2}. Its total space is the paths starting at \(b_*\), contracted by the homotopy after \eqref{mapsI:pf:eq:1.3}; its fibre is \(\Omega S^d\). The fibre is path connected because \(\pi_1(S^d)=0\): a based null-homotopy of a loop is a path to the constant loop, by \eqref{mapsI:pf:eq:1.1}. The sphere has its one-vertex, one-\(d\)-cell structure. With \(A_j=H_j(\Omega S^d;\mathbb Z)\), its cellular homology with these coefficients gives

\begin{equation}
E_2^{0,j}=A_j,\qquad E_2^{d,j}=A_j,
\label{mapsI:pf:eq:6.3}
\end{equation}

and no other columns. Only \(d_d:(d,j)\to(0,j+d-1)\) can be nonzero. No differential enters column \(d\), so its kernel must vanish, because the total space has no positive-degree homology. No differential leaves column zero, so its positive-degree quotients must vanish too. Thus

\begin{equation}
d_d:A_j\xrightarrow{\ \cong\ }A_{j+d-1}\quad(j\geq0).
\label{mapsI:pf:eq:6.4}
\end{equation}

The entries \(A_j\) with \(0<j<d-1\) have no incoming or outgoing differential, hence are zero. Connectivity gives \(A_0=\mathbb Z\). Induction using \eqref{mapsI:pf:eq:6.4} proves, in every degree,

\begin{equation}
H_j(\Omega S^5;\mathbb Z)=
\begin{cases}\mathbb Z,&4\mid j,\\0,&4\nmid j,\end{cases}
\qquad
H_j(\Omega S^6;\mathbb Z)=
\begin{cases}\mathbb Z,&5\mid j,\\0,&5\nmid j.\end{cases}
\label{mapsI:pf:eq:6.5}
\end{equation}

This determines the additive groups. No multiplicative structure or stable homotopy group is inferred from this calculation.

\section{The orientation generator and the path-fibration proof}
\label{mapsI:course:homotopy-section}

The homotopy-sphere conclusion is classical; the following argument retains
our original map and gives the distinct path-fibration proof of the course.
The complete elementary constructions used below are printed in
\cref{mapsI:cw:section,mapsI:pf:section}.  The first Hurewicz theorem and
homological Whitehead theorem are also stated in
\cite[Theorem~4.32 and Corollary~4.33]{Hatcher2002}; we use their integral,
simply connected forms.

\begin{theorem}[the original degree-one homotopy equivalence]
\label{mapsI:course:CG-S6-07-HOMOTOPY}
For the constructed threefold, with its complex orientation, there is a
map \(g:S^6\to X\) satisfying \(g_*[S^6]=[X]\), and this particular map
is a homotopy equivalence.
\end{theorem}

\begin{proof}
The fundamental-group calculation in
\cref{mapsI:top:prop:geometric-pi1} gives \(\pi_1(X)=0\), and
\cref{mapsI:top:thm:h2-vanishing,mapsI:top:lem:cohomological-closure}
give the complete integral homology of the sphere.  For \(2\le n\le5\),
assume inductively that \(\pi_i(X)=0\) for \(1\le i<n\).  The integral
Hurewicz map
\[
 h_n:\pi_n(X)\longrightarrow H_n(X;\mathbb Z),
 \qquad [a]\longmapsto a_*[S^n],
\]
is then an isomorphism, by \cref{mapsI:cw:hurewicz-theorem-result}.
The target vanishes, so the induction gives \(\pi_i(X)=0\) for
\(1\le i\le5\).  In degree six choose \([g]\) mapping to the positive
orientation generator \([X]\).  Thus
\begin{equation}
 g_*[S^6]=[X].
 \label{mapsI:course:degree-one}
\end{equation}
This map induces an integral homology isomorphism in every degree:
in degrees zero and six it preserves the generators, and all remaining
groups on both sides vanish.

Replace \(g\) by its actual path fibration
\[
 E_g=\{(s,\lambda):s\in S^6,\ \lambda(0)=g(s)\},
 \qquad p(s,\lambda)=\lambda(1).
\]
The path space has the compact-open topology.  The homotopy
\((s,\lambda)\mapsto(s,u\mapsto\lambda((1-t)u))\) strongly retracts
\(E_g\) onto the constant paths, identified with \(S^6\).
Under this retraction \(p\) represents \(g\).
For clarity, the lifting formula used in \cref{mapsI:pf:path-lifting}
is the following.  Given \(H:Z\times I\to X\) and an initial lift
\((s(z),\lambda_z)\), the path at time \(t\) is
\[
 \eta_{z,t}(u)=
 \begin{cases}
 \lambda_z((1+t)u),&0\le u\le (1+t)^{-1},\\
 H(z,(1+t)u-1),&(1+t)^{-1}\le u\le1.
 \end{cases}
\]
The pieces agree at the join, their two closed domains cover the parameter
space, and compact-open adjunction proves continuity.  The last endpoint
is \(H(z,t)\), while the time-zero path is exactly \(\lambda_z\).
This proves the homotopy lifting property.  Relative cell lifting gives
the exact homotopy sequence, as proved in
\cref{mapsI:pf:path-lifting}.

For the cellular filtration use the actual comparison constructed in
\cref{mapsI:cw:cw-manifolds}:
\[
 \alpha:T_X\longrightarrow X,\qquad
 \beta:X\longrightarrow T_X,\qquad
 \alpha\beta=1_X,\qquad \beta\alpha\simeq1_{T_X}.
\]
Put \(\bar g=\beta g\), let \(\bar E=E_{\bar g}\), and write
\(\bar p(s,\lambda)=\lambda(1)\).  There is an explicit comparison
\[
 J:\bar E\longrightarrow E_g,\qquad
 J(s,\lambda)=(s,\alpha\circ\lambda),
 \qquad pJ=\alpha\bar p.
\]
Its starting point is correct because
\(\alpha\beta g(s)=g(s)\), exactly.  Both total spaces retract onto the
same \(S^6\), and \(J\) commutes with these retractions.  The base
\(T_X\) is an actual simply connected CW complex with \(\pi_2(T_X)=0\).
The map \(\bar p\) is an integral homology isomorphism in every degree,
since \(\bar g=\beta g\) is one.  Let \(K\) be its fibre over a vertex.
It is path connected: a path between two fibre points in the connected
total space projects to a null-homotopic loop, and lifting its contraction
relative to the two endpoints changes that path to one in the fibre.
The exact sequence, together with \(\pi_2(T_X)=0\) and
\(\pi_1(\bar E)=\pi_1(S^6)=0\), then gives \(\pi_1(K)=0\).

Suppose a positive homotopy group of \(K\) is nonzero and choose the least
such degree \(n\ge2\).  Hurewicz gives
\[
 H_i(K;\mathbb Z)=0\quad(0<i<n),\qquad
 H_n(K;\mathbb Z)\cong\pi_n(K)\ne0.
\]
The integral homology spectral sequence for \(\bar p\), constructed
with its actual edge maps in \cref{mapsI:pf:filtered-chains,mapsI:pf:serre-homology},
has constant coefficient systems because the base is simply connected.
In total degree \(n\), the only possibly nonzero terms are
\[
 E_2^{0,n}=H_n(K),\qquad E_2^{n,0}=H_n(T_X).
\]
No differential leaves column zero.  A differential entering \((0,n)\)
has source \((r,n-r+1)\).  The sources for \(2\le r\le n\) vanish by
the displayed fibre-homology vanishings, and for \(r>n+1\) they have
negative second degree.  The sole possible incoming map is therefore
\[
 d_{n+1}:E_{n+1}^{n+1,0}\longrightarrow E_{n+1}^{0,n}.
\]
The source is still \(H_{n+1}(T_X)\): there are no incoming maps to
the bottom row, and all earlier outgoing targets have fibre degree
between one and \(n-1\).  The actual edge image
\(\operatorname{im}(\bar p_*:H_{n+1}(\bar E)\to H_{n+1}(T_X))\)
is \(E_\infty^{n+1,0}=\ker d_{n+1}\).  Surjectivity of \(\bar p_*\)
makes this kernel the whole source, so \(d_{n+1}=0\).
Consequently
\[
 H_n(K)=E_\infty^{0,n}=F_0H_n(\bar E).
\]
The other actual edge map \(\bar p_*:H_n(\bar E)\to H_n(T_X)\)
kills this first filtration subgroup: its kernel is
\(F_{n-1}H_n(\bar E)\), which contains \(F_0\) since \(n\ge2\).
Injectivity of this edge map therefore forces \(H_n(K)=0\), a
contradiction.  All homotopy groups of \(K\) vanish.

The exact homotopy sequence makes \(\bar g\), and hence the original
\(g\), an isomorphism on every homotopy group.  The complete CW-domain
compression argument in \cref{mapsI:pf:cw-inverses} supplies an actual
inverse \(\bar h:T_X\to S^6\) to \(\bar g\), with both inverse
homotopies.  Set
\[
 h_X=\bar h\beta:X\longrightarrow S^6.
\]
Then \(h_Xg=\bar h\bar g\simeq1_{S^6}\).  For the other composition,
use \(g=\alpha\bar g\), then \(\bar g\bar h\simeq1_{T_X}\), then
the exact equality \(\alpha\beta=1_X\), to obtain
\[
 gh_X=\alpha\bar g\bar h\beta\simeq\alpha\beta=1_X.
\]
Both homotopies concern the original spaces and the original degree-one
map.  The orientation generator in \eqref{mapsI:course:degree-one}
has not been rescaled.
\end{proof}
\begin{thebibliography}{1}
\bibitem{Hatcher2002} Allen Hatcher. Algebraic Topology. Cambridge University Press, 2002.
\end{thebibliography}
\end{document}
